%%%%%%%%%%%%%%%%%%%%%%%%%%%%%%%%%%%%%%%%%%%%%%%%%%%%%%%%%%%%%%%%%%%%%%%%%%%%%
%                                                                           %
%          LaTeX File for Doctor (Master) Thesis of ECNU                    %
%            华东师范大学博士（硕士）论文模板 ____lizb                        %
%                                                                           %
%%%%%%%%%%%%%%%%%%%%%%%%%%%%%%%%%%%%%%%%%%%%%%%%%%%%%%%%%%%%%%%%%%%%%%%%%%%%%
%!TEX program = xelatex
%!BIB program = bibtex
\documentclass[12pt,openany,a4paper,fancyhdr,oneside]{ctexbook}
%\documentclass[12pt,openright,a4paper,fancyhdr,twoside]{ctexbook}
%draft 选项可以使插入的图形只显示外框，以加快预览速度。
%\documentclass[11pt,a4paper,openany,draft]{book}

\usepackage{multirow}
\usepackage{listings}
\usepackage{xcolor}

\usepackage[CJKbookmarks,linkcolor=black,citecolor=black]{hyperref}
\usepackage{pdfpages}
\usepackage{shortvrb,indentfirst,ulem,makeidx}
\usepackage{fancyhdr}
\usepackage{graphicx}
\usepackage{indentfirst,latexsym,amsthm,colortbl,subfigure,clrscode}
\usepackage{algorithm}
% \usepackage{algorithmicx}
% \usepackage{algpseudocode}   % 如果后面出问题，注释掉这个
\usepackage{algorithmic}    % 原本是有这个，和上面的冲突，注释了
\usepackage{bm}                     % 处理数学公式中的黑斜体的宏包
\usepackage{amsmath}                % AMSLaTeX宏包 用来排出更加漂亮的公式
\usepackage{amssymb}                % AMSLaTeX宏包 用来排出更加漂亮的公式

\usepackage{mathrsfs}
\usepackage[subnum]{cases}
\usepackage[numbers,sort&compress]{natbib}
%\usepackage[super,square,numbers,sort&compress]{natbib}
\usepackage{hypernat}
\usepackage{enumerate}
\usepackage{geometry}

\usepackage{times}
\usepackage{fontspec}
\usepackage{libertine}

\usepackage{caption}
\usepackage{titletoc}
%下面是我自己加的包
\usepackage{multirow}
\usepackage{bigdelim}
\usepackage{caption}
\usepackage{amsmath}
\usepackage{graphicx,color,xcolor,colortbl}
\usepackage{epstopdf,epsfig}
\usepackage{amssymb}
\usepackage{tabularx}
\DeclareMathOperator*{\argmax}{argmax}
\DeclareMathOperator*{\argmin}{argmin}
% \usepackage{color}
% \usepackage[dvipsnames]{xcolor}
\hyphenation{op-tical net-works semi-conduc-tor}


\makeindex
\pagestyle{fancy}

\fancyhead[RO,LE]{\wuhao{华东师范大学研究生硕士学位论文}}
\fancyhead[LO]{\wuhao \leftmark} 
\fancyhead[RE]{\wuhao \leftmark}

\renewcommand{\headrulewidth}{0.4pt}
\fancyfoot[CO,CE]{\thepage}

\renewcommand{\algorithmicrequire}{\textbf{Input:}}
\renewcommand{\algorithmicensure}{\textbf{Output:}}
\renewcommand{\algorithmiccomment}[1]{// #1}


%                    根据自己正文需要做的一些定义                 %
%==================================================================
\def\diag{{\rm diag}}
\def\rank{{\rm rank}}
\def\RR{{\cal R}}
\def\NN{{\cal N}}
\def\R{{\mathbb R}}
\def\C{{\mathbb C}}
\let\dis=\displaystyle

\def\p{\partial}
\def\f{\frac}
\def\mr{\mathrm}
\def\mb{\mathbf}
\def\mc{\mathcal}
\def\b{\begin}
\def\e{\end}

\newtheorem{thm1}{Theorem}[part]
\newtheorem{thm2}{Theorem}[section]
\newtheorem{thm3}{Theorem}[subsection]
\newtheorem{them}[thm2]{定理}
\newtheorem{theorem}[thm2]{定理}
\newtheorem{defn}[thm2]{定义}
\newtheorem{define}[thm2]{定义}
\newtheorem{ex}[thm2]{例}
\newtheorem{exs}[thm2]{例}
\newtheorem{example}[thm2]{例}
\newtheorem{prop}[thm2]{命题}
\newtheorem{lemma}[thm2]{引理}
\newtheorem{cor}[thm2]{推论}
\newtheorem{remark}[thm2]{注释}
\newtheorem{notation}[thm2]{记号}
\newtheorem{abbre}[thm2]{缩写}
% \newtheorem{algorithm}[thm2]{算法}
\newtheorem{problem}[thm2]{问题}
\newcommand{\gameauth}{\mathsf{Game}^\mathsf{Auth}}
\newcommand{\gametul}{\mathsf{Game}^\mathsf{TUL}}
\newcommand{\gamepriv}{\mathsf{Game}^\mathsf{PRIV}}
\newcommand{\advaauth}{\mathsf{Adv}_{\mc A}^{\mathsf{Auth}}}
\newcommand{\advatul}{\mathsf{Adv}_{\mc A}^{\mathsf{TUL}}}
\newcommand{\advapriv}{\mathsf{Adv}_{\mc A}^{\mathsf{PRIV}}}
\newcommand{\tsf}{\textsf}
\newtheorem{analysis}{Analysis}
\newcommand{\be}{\begin{enumerate}}
\newcommand{\ed}{\end{enumerate}}

\newcommand{\tabincell}[2]{\begin{tabular}{@{}#1@{}}#2\end{tabular}}
\newcommand{\mbb}{\mathbb}
\newcommand{\yihao}{\fontsize{26pt}{36pt}\selectfont}           % 一号, 1.4 倍行距
\newcommand{\erhao}{\fontsize{22pt}{28pt}\selectfont}          % 二号, 1.25倍行距
\newcommand{\xiaoer}{\fontsize{18pt}{18pt}\selectfont}          % 小二, 单倍行距
\newcommand{\sanhao}{\fontsize{16pt}{24pt}\selectfont}        % 三号, 1.5倍行距
\newcommand{\xiaosan}{\fontsize{15pt}{22pt}\selectfont}        % 小三, 1.5倍行距
\newcommand{\sihao}{\fontsize{14pt}{21pt}\selectfont}            % 四号, 1.5 倍行距
\newcommand{\banxiaosi}{\fontsize{13pt}{19.5pt}\selectfont}    % 半小四, 1.5倍行距
\newcommand{\xiaosi}{\fontsize{12pt}{18pt}\selectfont}            % 小四, 1.5倍行距
\newcommand{\dawuhao}{\fontsize{11pt}{11pt}\selectfont}       % 大五号, 单倍行距
\newcommand{\wuhao}{\fontsize{10.5pt}{15.75pt}\selectfont}    % 五号, 单倍行距



\lstset
{
	basicstyle=\ttfamily,
	% numbers=left,
	% numberstyle=\tiny,
	keywordstyle=\color[RGB]{0, 0, 255},
	commentstyle=\color[RGB]{0, 128, 0},
	frame=shadowbox,
	rulesepcolor=\color{red!20!green!20!blue!20},
	showspaces=false,
	showstringspaces=false,
	extendedchars=false,
	showtabs=false,
	tabsize=4,
	xleftmargin=0.5em,
	xrightmargin=0.5em,
	% aboveskip=1em,
	escapeinside=``
}


%============================ 可以自定义文字块 ================================%

\newcommand{\aaa}{这是我给你们的一个示例}
\newcommand{\bbb}{\aaa \aaa \aaa}
\newcommand{\ccc}{\bbb \bbb \bbb \bbb \bbb

\bbb \bbb \bbb \bbb \bbb }
\newcommand{\abc}{abcdefg1234567890}
\newcommand{\upabc}{ABCDEFGHIJK}
%%% ----------------------------------------------------------------------



%============================= 版芯控制 ================================%
\setlength{\oddsidemargin}{0.57cm}
\setlength{\evensidemargin}{\oddsidemargin}
\voffset-6mm \textwidth=150mm \textheight=230mm \headwidth=150mm
%\rightmargin=35mm
%                                                                       %


%============================= 页面设置 ================================%
%-------------------- 定义页眉和页脚 使用fancyhdr 宏包 -----------------%
% 定义页眉与正文间双隔线
\newcommand{\makeheadrule}{%
\makebox[0pt][l]{\rule[.7\baselineskip]{\headwidth}{0.4pt}}%
\rule[0.85\baselineskip]{\headwidth}{0.4pt} \vskip-.8\baselineskip}
\makeatletter
\renewcommand{\headrule}{%
{\if@fancyplain\let\headrulewidth\plainheadrulewidth\fi
\makeheadrule}} \makeatother

\newcommand{\adots}{\mathinner{\mkern 2mu%
\raisebox{0.1em}{.}\mkern 2mu\raisebox{0.4em}{.}%
\mkern2mu\raisebox{0.7em}{.}\mkern 1mu}}

%\setmainfont{Times New Roman}
\dottedcontents{chapter}[1.5cm]{\xiaosi\heiti}{3.8em}{9.5pt}
\dottedcontents{section}[1.5cm]{\xiaosi\heiti}{2.8em}{9.5pt}


%=============================== 正文部分 ================================%
\begin{document}

\pagestyle{empty}
\setlength{\baselineskip}{25pt}  %%正文设为25磅行间距
\vspace{-2.0cm}
\noindent{{\zihao{4} {\large 2020} 届研究生硕士学位论文}}\\
\vspace{-0.8cm}
\begin{flushleft}
\hspace{-0.5cm}
\renewcommand\arraystretch{1.5}
\begin{tabular}{l}
\noindent{{\zihao{4} 分类号：\underline{\qquad\qquad\qquad\qquad\qquad\qquad}}}  \\
\noindent{{\zihao{4} 密~~~~级：\underline{\qquad\qquad\qquad\qquad\qquad\qquad}}}\\
\end{tabular}
\hskip 1.1cm
\renewcommand\arraystretch{1.5}
\begin{tabular}{l}
\noindent{{\zihao{4} 学校代码：\underline{10269~~~\qquad}}}\\
% \noindent{{\zihao{4} 学~~~~~~~~号：\underline{}}}\\
\noindent{{\zihao{4} 学~~~~~~~~号：\underline{51170000000}}}\\
\end{tabular}
\end{flushleft}


\vskip 1.0cm

\begin{center}
	\hskip 0.5cm
	\scalebox{1.0}{\includegraphics[width=2.7cm]{fig/ecnulogo.png}}
	\scalebox{1.0}{\includegraphics[width=11.5cm,height=2.7cm]{fig/ecnulabel.png}}
	\vskip 0.5cm
	{\textbf{{\xiaoer East China Normal University}}}\\ \vskip 0.2cm
	{\textbf{\erhao 硕~士~学~位~论~文}}\\ \vskip 0.2cm
	{\textbf{{\xiaoer MASTER'S DISSERTATION}}}\\\end{center}





\vskip 1.0cm

%基于分数阶微分和特征的一致性图像配准研究
\begin{center}
{\erhao \bf 论文题目：\underline{~~你的论文标题在这里~~}}\\
{\erhao \bf ~~~~~~~~~~~~~~~~~~~~\underline{~~你的论文标题在这里~~}}
%{\erhao \bf ~~~~~~~~~~~~~~~~~~~~\underline{~~~~~~~~~~~~~~~~}}
\end{center}
\vskip 1.0cm
\begin{center}

\renewcommand\arraystretch{1.5}
	\begin{tabular}{l}
{\sihao \bf 院\qquad\ \ \ 系:}\\
{\sihao \bf 专~业~名~称:}\\
{\sihao \bf 研~究~方~向:}\\
{\sihao \bf 指~导~教~师:}\\
{\sihao \bf 学位申请人:}
\end{tabular}
\begin{tabular}c
{\sihao \bf  ~~计算机科学与技术学院}               \\
\hline {\sihao \bf ~~计算机科学与技术}              \\
\hline {\sihao \bf ~~图像处理与模式识别~~}\\
% \hline {\sihao \bf ~~ \ \ }  \\
% \hline{\sihao \bf  ~~}      \\
\hline {\sihao \bf ~~XX~~教授\ \ }  \\
\hline{\sihao \bf  ~~XX}      \\
\hline
\end{tabular}


\end{center}

\vskip 1.95cm
\begin{center}
{\sihao 2020年X月X日}
\end{center}

%\clearpage\ \newpage
\newpage

\pagestyle{empty}

%\noindent{\large Dissertation for Master Degree in 2020}
%\hskip 1cm {\large University Code: 10269}\\
%\hspace *{\fill}{\large Student Number: 51174506034}
\noindent{\large 2020 MASTER’S  DISSERTATION}
\hskip 1.4cm {\large School Code: 10269}\\
\hspace*{\fill}{\large Student Number: 51170000000}

\vskip 2cm

\begin{center}
{\Huge $\mathbb{EAST}\,\mathbb{CHINA}\,\mathbb{NORMAL}\,
\mathbb{UNIVERSITY}$}
\end{center}

\vskip 3cm

%\begin{center}
%{\erhao\bf  Title: \underline{The Combination of Side Window Fractional}}\\
%{\erhao\bf  \underline{~~~~Order Derivate and Feature with Inertial~~~}}\\
%{\erhao\bf  \underline{Constraint for Non-rigid Image Registration}}
%\end{center}
\begin{center}
{\erhao\bf  Title: \underline{Research on XXXXX XXXXXX }}\\
{\erhao\bf  \underline{~~~XXXXXXX XXXXXXX XXXXX~~~}}\\
{\erhao\bf  \underline{~~~~~~XXXX XXXXXXXX XXXXXXX~~~~~~~}}
\end{center}

\vskip 2cm {\large
\begin{center}
\begin{tabular}{l}
Department:\\
Major:\\
Research Area:\\
Supervisor:\\
Candidate:
\end{tabular}
\begin{tabular}c
~~~School of Computer Science and Technology \\
\hline ~~~Computer Science and Technology  \\
\hline ~~~Image Processing and Pattern Recognition\\
% \hline ~~~  \\
% \hline ~~~ \\
\hline ~~~XXXX~ Professor  \\
\hline ~~~XXXX  \\
\hline
\end{tabular}
\end{center}}

\vskip 30mm

\begin{center}
{\Large Month, 2020}
\end{center}

%\clearpage\ \newpage
\newpage
\pagestyle{empty}
% \includepdf{fig/signed}        原创性申明的签署页，作为pdf文件放在fig文件夹下
\centerline{\bf\Large 华东师范大学学位论文原创性声明}

\vskip 0.8cm

\normalsize \indent
郑重声明：本人呈交的学位论文《XXXX（你的论文题目）》，是在华东师范大学攻读硕士/博士（请勾选）学位期间，在导师的指导下进行的研究工作及取得的研究成果。除文中已经注明引用的内容外，本论文不包含其他个人已经发表或撰写过的研究成果。对本文的研究做出重要贡献的个人和集体，均已在文中作了明确说明并表示谢意。
$$\\  $$

\qquad\qquad{作者签名}:$\underline{\qquad\qquad\qquad\qquad }$
\qquad \qquad\qquad \mbox {日期}:\qquad 年 \qquad  月 \qquad  日


\vskip 0.8cm

\centerline{\bf\Large 华东师范大学学位论文著作权使用声明}

\vskip 0.8cm

《XXXX（你的论文题目）》系本人在华东师范大学攻读学位期间在导师指导下完成的硕士/博士（请勾选）学位论文，本论文的研究成果归华东师范大学所有。本人同意华东师范大学根据相关规定保留和使用此学位论文，并向主管部门和相关机构如国家图书馆、中信所和“知网”送交学位论文的印刷版和电子版；允许学位论文进入华东师范大学图书馆及数据库被查阅、借阅；同意学校将学位论文加入全国博士、硕士学位论文共建单位数据库进行检索，将学位论文的标题和摘要汇编出版，采用影印、缩印或者其它方式合理复制学位论文。

本学位论文属于（请勾选）

（~~）1.经华东师范大学相关部门审查核定的“内部”或“涉密”学位论文*，
于年~~~~~~月~~~~~~日解密，解密后适用上述授权。

（~~）2.不保密，适用上述授权。
$$\\ $$
\qquad\qquad \mbox{导师签名}:$\underline{\qquad\qquad\qquad\qquad}$
\qquad\qquad \mbox {本人签名}:$\underline{\qquad\qquad\qquad\qquad }$

\vskip 0.8cm

$\rightline{ \qquad 年 \qquad  月 \qquad  日 \qquad\qquad}$

\vskip 0.8cm

* “涉密”学位论文应是已经华东师范大学学位评定委员会办公室或保密委员会审定过的学位论文（需附获批的《华东师范大学研究生申请学位论文“涉密”审批表》方为有效），未经上述部门审定的学位论文均为公开学位论文。此声明栏不填写的，默认为公开学位论文，均适用上述授权）。

%\clearpage\ \newpage
\newpage
\pagestyle{empty}
$$\\ \\ \\ $$

\centerline{\bf\Large $\underline{\mbox{\kaishu {}}}\,\,
	$硕士学位论文答辩委员会成员名单}

\vskip 10mm

\begin{center}
	{\large
		\renewcommand{\arraystretch}{1.4}
		\begin{tabular}{|c|c|c|c|} \hline
			{\heiti ~~~~~姓~名~~~~~} &{\heiti ~~~~~职~称~~~~~} 
			& \hspace{6em}{\heiti 单~位}\hspace{6em} & {\heiti ~~~~备~注~~~~}\\\hline
			{\kaishu XXX}    &   {\kaishu 教授}    &  {\kaishu 华东师范大学}  & {\kaishu 主席}  \\ \hline
			{\kaishu XXX}    &   {\kaishu 教授}   &  {\kaishu 华东师范大学} &       \\ \hline
			{\kaishu XXX}    &   {\kaishu 教授}    &  {\kaishu 华东师范大学}&       \\ \hline
			&           &                   &       \\ \hline
			&           &                   &       \\ \hline
		\end{tabular}
	}
\end{center}


%\clearpage\ \newpage

%\newpage
\pagenumbering{Roman}  %摘要页的页码是罗马字符 roman是小写  Roman是大写
\pagestyle{plain}

\addcontentsline{toc}{chapter}{摘要}   %在目录页添加摘要及页码，如不需要可直接注释
\vspace{-2.5cm}
\chapter*{\zihao{2}\heiti{摘~~~~要}}
%\vskip 1cm
%\vspace{-1cm}

论文摘要论文摘要论文摘要论文摘要论文摘要论文摘要论文摘要论文摘要论文摘要论文摘要论文摘要论文摘要论文摘要论文摘要论文摘要论文摘要论文摘要论文摘要论文摘要论文摘要论文摘要论文摘要论文摘要论文摘要论文摘要论文摘要论文摘要论文摘要论文摘要论文摘要论文摘要论文摘要论文摘要论文摘要论文摘要论文摘要论文摘要论文摘要论文摘要论文摘要论文摘要论文摘要论文摘要论文摘要论文摘要论文摘要论文摘要论文摘要论文摘要论文摘要论文摘要论文摘要论文摘要论文摘要论文摘要论文摘要论文摘要论文摘要论文摘要论文摘要论文摘要论文摘要

本文的主要工作包括：
\begin{enumerate}
	\item
	\textbf{创新点1创新点1创新点1}论文摘要论文摘要论文摘要论文摘要论文摘要论文摘要论文摘要论文摘要论文摘要论文摘要论文摘要论文摘要论文摘要论文摘要论文摘要论文摘要论文摘要论文摘要论文摘要论文摘要论文摘要论文摘要论文摘要论文摘要论文摘要论文摘要论文摘要论文摘要论文摘要论文摘要论文摘要
	
	\item
	\textbf{创新点2创新点2创新点2}论文摘要论文摘要论文摘要论文摘要论文摘要论文摘要论文摘要论文摘要论文摘要论文摘要论文摘要论文摘要论文摘要论文摘要论文摘要论文摘要论文摘要论文摘要论文摘要论文摘要论文摘要论文摘要论文摘要论文摘要论文摘要论文摘要论文摘要论文摘要论文摘要论文摘要论文摘要
		
	\item
	\textbf{创新点3创新点3创新点3}论文摘要论文摘要论文摘要论文摘要论文摘要论文摘要论文摘要论文摘要论文摘要论文摘要论文摘要论文摘要论文摘要论文摘要论文摘要论文摘要论文摘要论文摘要论文摘要论文摘要论文摘要论文摘要论文摘要论文摘要论文摘要论文摘要论文摘要论文摘要论文摘要论文摘要论文摘要
\end{enumerate}

论文摘要论文摘要论文摘要论文摘要论文摘要论文摘要论文摘要论文摘要论文摘要论文摘要论文摘要论文摘要论文摘要论文摘要论文摘要论文摘要论文摘要论文摘要论文摘要论文摘要论文摘要论文摘要论文摘要论文摘要论文摘要论文摘要论文摘要论文摘要论文摘要论文摘要论文摘要

\hspace{-0.5cm}
\sihao{\heiti{关键词：}} \xiaosi{关键词，keyword，关键词，keyword，关键词，keyword}


\addcontentsline{toc}{chapter}{ABSTRACT}  %在目录页添加abstract及页码，如不需要可直接注释
\newpage
\vspace{-1cm}
\chapter*{\zihao{-2}\heiti{ABSTRACT}}
%\vspace{-0.5cm}
Your abstract. Your abstract. Your abstract. Your abstract. Your abstract. Your abstract. Your abstract. Your abstract. Your abstract. Your abstract. Your abstract. Your abstract. Your abstract. Your abstract. Your abstract. Your abstract. Your abstract. Your abstract. Your abstract. Your abstract. Your abstract. Your abstract. Your abstract. Your abstract. Your abstract. Your abstract.

The main work of this paper includes:
\begin{enumerate}
\item 
\textbf{First of your idea.First of your idea.First of your idea.First of your idea.First of your idea.}
Your abstract. Your abstract. Your abstract. Your abstract. Your abstract. Your abstract. Your abstract. Your abstract. Your abstract. Your abstract. Your abstract. Your abstract. Your abstract. Your abstract. Your abstract. Your abstract. Your abstract. Your abstract. Your abstract. Your abstract. Your abstract. Your abstract. Your abstract. Your abstract. Your abstract. Your abstract.
	
\item
\textbf{Second of your idea.Second of your idea.Second of your idea.Second of your idea.Second of your idea.} Your abstract. Your abstract. Your abstract. Your abstract. Your abstract. Your abstract. Your abstract. Your abstract. Your abstract. Your abstract. Your abstract. Your abstract. Your abstract. Your abstract. Your abstract. Your abstract. Your abstract. Your abstract. Your abstract. Your abstract. Your abstract. Your abstract. Your abstract. Your abstract. Your abstract. Your abstract.

\item
\textbf{Third of your idea. Third of your idea.Third of your idea. Third of your idea. Third of your idea.} Your abstract. Your abstract. Your abstract. Your abstract. Your abstract. Your abstract. Your abstract. Your abstract. Your abstract. Your abstract. Your abstract. Your abstract. Your abstract. Your abstract. Your abstract. Your abstract. Your abstract. Your abstract. Your abstract. Your abstract. Your abstract. Your abstract. Your abstract. Your abstract. Your abstract. Your abstract.
\end{enumerate}

Your abstract. Your abstract. Your abstract. Your abstract. Your abstract. Your abstract. Your abstract. Your abstract. Your abstract. Your abstract. Your abstract. Your abstract. Your abstract. Your abstract. Your abstract. Your abstract. Your abstract. Your abstract. Your abstract. Your abstract. Your abstract. Your abstract. Your abstract. Your abstract. Your abstract. Your abstract.

%空行不能省
%\hspace{-0.5cm}
{\sihao{\textbf{Keywords:}}} \textit{keywords; keywords; keywords; keywords}




































\setcounter{tocdepth}{2}

\addcontentsline{toc}{chapter}{目录}  %在目录页添加目录及页码，如不需要可直接注释
\tableofcontents

\addcontentsline{toc}{chapter}{插图}  %在目录页添加插图及页码，如不需要可直接注释
\listoffigures

\addcontentsline{toc}{chapter}{表格}  %在目录页添加表格及页码，如不需要可直接注释
\listoftables

\newpage
\pagenumbering{arabic}
\pagestyle{fancy}


\CTEXsetup[format+={\zihao{3}\heiti}]{chapter}
\CTEXsetup[format+={\raggedright\zihao{4}\heiti}]{section}
\CTEXsetup[format+={\zihao{-4}\heiti}]{subsection}


\setlength{\baselineskip}{25pt}  %%正文设为25磅行间距

\chapter{绪论}
\label{ch1}
\section{研究背景与意义}
研究背景与意义研究背景与意义研究背景与意义研究背景与意义研究背景与意义研究背景与意义研究背景与意义研究背景与意义研究背景与意义研究背景与意义研究背景与意义研究背景与意义研究背景与意义研究背景与意义研究背景与意义研究背景与意义研究背景与意义研究背景与意义\cite{ReferenceName}。如图\ref{fig:traditional_framework}所示。

\begin{figure}[H]\vspace{0pt}
	\centering
	\includegraphics[width=0.95\linewidth]{fig/timg}
	\caption{图片的题注}
	\label{fig:traditional_framework}\vspace{0pt}
\end{figure}

\subsection{小标题1}
小标题1小标题1小标题1小标题1小标题1小标题1小标题1小标题1小标题1小标题1小标题1小标题1小标题1小标题1小标题1小标题1小标题1小标题1小标题1小标题1小标题1小标题1小标题1小标题1小标题1小标题1小标题1小标题1小标题1小标题1小标题1小标题1小标题1小标题1小标题1小标题1小标题1小标题1小标题1小标题1小标题1

\subsection{小标题2}
小标题2小标题2小标题2小标题2小标题2小标题2小标题2小标题2小标题2小标题2小标题2小标题2小标题2小标题2小标题2小标题2小标题2小标题2小标题2小标题2小标题2小标题2小标题2小标题2小标题2小标题2小标题2小标题2小标题2小标题2小标题2小标题2小标题2小标题2小标题2小标题2小标题2小标题2小标题2小标题2小标题2小标题2小标题2小标题2小标题2小标题2小标题2小标题2小标题2小标题2小标题2小标题2小标题2小标题2小标题2小标题2小标题2小标题2小标题2小标题2。如表\ref{tab:table_eg_1}所示

\begin{table}[h!]
	\centering
	\caption{表格的题注}
	\label{tab:table_eg_1}
	\begin{tabular}{l|ll}
		\hline
		xx & xx & xx \\ \hline
		xx & 11 & 11 \\
		xx & 11 & 11 \\ \hline
	\end{tabular}
\end{table}

\section{本文主要工作}
本文主要工作本文主要工作本文主要工作本文主要工作本文主要工作本文主要工作本文主要工作本文主要工作本文主要工作本文主要工作本文主要工作本文主要工作本文主要工作本文主要工作本文主要工作本文主要工作本文主要工作本文主要工作本文主要工作本文主要工作本文主要工作本文主要工作本文主要工作本文主要工作


%\input{C1-CHAP2.tex}
%\input{C1-CHAP3.tex}
%\input{C1-CHAP4.tex}
%\input{C1-CHAP5.tex}
%\input{C1-CHAP6.tex}
%\input{C8-CHAP8.tex}
%\appendix

\chapter{附录}
\vspace{-1cm}







\pagestyle{plain}
\clearpage
\phantomsection
\addcontentsline{toc}{chapter}{参考文献}
% \begin{thebibliography}{zz}




% \bibitem{2000Joux}
% Joux, Antoine. "A one round protocol for tripartite Diffie–Hellman."Algorithmic number theory. Springer Berlin Heidelberg, 2000. 385-393.

% \bibitem{2001Boneh}

% Boneh, Dan, and Matt Franklin. "Identity-based encryption from the Weil pairing." Advances in Cryptology—CRYPTO 2001. Springer Berlin Heidelberg, 2001.

% \bibitem{1984S}
% Identity-Based Cryptosystems and Signature Schemes

% \bibitem{2003J}
% J.C.Cha,J.H.Cheon.An identity-based signature from gap Diffie-Hellman groups.In Y Desmedt,editor,Public Key Cryptography-PKC 2003,LNCS 2567,Berlin:Springer-Verlag,2003:18-30



% \end{thebibliography}


\bibliographystyle{bib/GBT7714-2005}
\bibliography{bib/tex}

%\pagestyle{plain}\clearpage
\pagestyle{plain}
\clearpage
\phantomsection
\addcontentsline{toc}{chapter}{致谢}
{\fangsong
	\chapter*{致\qquad 谢}\vskip 2mm
	\vspace{-1cm}
	\large{致谢致谢致谢致谢致谢致谢致谢致谢致谢致谢致谢致谢致谢致谢致谢致谢致谢致谢致谢致谢致谢致谢致谢致谢致谢致谢致谢致谢致谢致谢致谢致谢致谢致谢致谢致谢致谢致谢致谢致谢致谢致谢致谢致谢致谢致谢致谢致谢



%\vspace{0.2cm} \hspace{9.8cm}  
%徐~~诚
%
%\hspace{9cm}  二零二零年三月
  

	}
	
	\vspace{0.2cm}
	

}

\pagestyle{plain}
\clearpage
\phantomsection
\addcontentsline{toc}{chapter}{发表论文和科研情况}
\chapter*{\large 攻读硕士学位期间参与的项目以及学术成果}
\vskip 2mm
\vspace{-1cm}
\renewcommand{\labelenumi}{[\arabic{enumi}]}


{\heiti $\blacksquare$ 已完成学术论文}\vskip 3mm
\begin{enumerate}
	\item \textbf{XXXXXXX}, XXXXXXXX XXXXXXXXXXXX XXXXXXXXXXXXXXX XXXXXX XXXXXXXXXX XXXXXX
	
%	\item Gabor Feature Based LogDemons with Inertial Constraint for Nonrigid Image Registration. IEEE Transactions on Image Processing. (Major Revision)
\end{enumerate} 




\printindex
\end{document}

